\subsection{对偶化模}

定义 1.—— 设 A 为 Noether 环。称 Λ-模 Ω 为对偶化的，如果它是有限生成的，并且对 Λ 的每个极大理想 m，A/m-向量空间 Ext \(_{A}^{i}\) (A/m, Ω) 在 i ≠ ht(m) 时为零，而在 i = ht(m) 时维数为 1。

对 A 的每个极大理想 m 与每个整数 i，A/m-向量空间 \(\operatorname{Ext}_{\mathrm{A}}^{i}(\mathrm{A}/\mathfrak{m},\Omega)\) 典范地同构于 \(\operatorname{Ext}_{\mathrm{A}_{\mathfrak{m}}}^{i}(\mathrm{A}/\mathfrak{m},\Omega_{\mathfrak{m}})\)。因此，为使有限生成 A-模 \(\Omega\) 是对偶化的，必须且只须 \(A_{m}\)-模 \(\Omega_{m}\) 对 A 的每个极大理想 m 都是对偶化的。

例.—— 1) 若环 A 是局部 Artin 环，则对偶化 A-模就是有限生成的内射 A-模 \(\Omega\)，它使得 \(\mathrm{Hom}_{\mathrm{A}}(\kappa_{\mathrm{A}},\Omega)\)（记之为 \( E \)）的维数等于 1（\(\S\) 3, 第 3 节，命题 6），亦即 Matlis A-模（\(\S\) 8, 第 3 节）。

\begin{enumerate}
\def\labelenumi{\arabic{enumi})}
\setcounter{enumi}{1}
\tightlist
\item
  为使 Noether 环 A 是 Gorenstein 环，必须且只须 A-模 A 是对偶化的（§ 3, 第 7 节，命题 11）。特别地，当 A 正则时，A-模 A 是对偶化的。
\end{enumerate}

注.—— 1) 设 A 为 Noether 局部环，\(\Omega\) 为有限生成 A-模。剩余域 \(\kappa_{\widehat{A}}\) 等同于 \(\kappa_{A}\)，而 \(\widehat{A}\)-模 \(\widehat{\Omega}\) 等同于 \(\widehat{A} \otimes_{A} \Omega\)（III, § 3, 第 4 节，定理 3）。于是由 A, X, 第 111 页，命题 10 可知，\(\kappa_{A}\)-向量空间 \(\mathrm{Ext}_{\mathbf{A}}^{i}(\kappa_{\mathbf{A}}, \Omega)\) 典范地同构于 \(\mathrm{Ext}_{\widehat{A}}^{i}(\kappa_{\widehat{A}}, \widehat{\Omega})\)。因此，为使 \(A\)-模 \(\Omega\) 是对偶化的，必须且只须 \(\widehat{A}\)-模 \(\widehat{\Omega}\) 是对偶化的。

\begin{enumerate}
\def\labelenumi{\arabic{enumi})}
\setcounter{enumi}{1}
\tightlist
\item
  设 \(\Omega\) 为对偶化 A-模；对每个秩 1 的投射 A-模 L，A-模 \(\Omega\otimes_{\mathrm{A}}\mathrm{L}\) 是对偶化的（A, X, 第 108 页，命题 7, b)）。我们将在下面（\(n^{\circ}\) 4, 命题 6）看到，每个对偶化 A-模都同构于这种形状的模。
\end{enumerate}

命题 1.—— 设 A 为 Noether 环，Ω 为对偶化 A-模。

\begin{enumerate}
\def\labelenumi{\alph{enumi})}
\tightlist
\item
  则 A 是 Macaulay 环，且 A-模 Ω 是 Macaulay 模。
\end{enumerate}

\[
\text { b) } O n a \mathrm{di} _ {\mathrm{A}} (\Omega) = \dim (\Omega) = \dim (\mathrm{A}).
\]

先设环 A 是局部环，并用 \(d\) 表示它的维数。§ 3, 第 3 节的命题 6 蕴含 \(\mathrm{di}_{\Lambda}(\Omega) = d\)，从而由 § 3, 第 6 节的命题 9 得 \(\mathrm{prof}(\mathrm{A}) = d\)，故 A 是 Macaulay 环。此外，按深度的定义，有 \(\mathrm{prof}(\Omega) = d\)；由于有 \(\mathrm{prof}(\Omega) \leqslant \dim(\Omega) \leqslant d\)，由此即得此情形下的命题。

在一般情形，\(A_{m}\)-模 \(\Omega_{m}\) 对 A 的每个极大理想 m 都是对偶化的，从而 \(A_{m}\) 是 Macaulay 环，而 \(\Omega_{m}\) 是 \(\Lambda_{m}\)-模 a)。此外，对每个极大理想 m，有 \(\mathrm{di}_{\mathrm{A}_{\mathfrak{m}}(\Omega_{\mathfrak{m}})} = \dim(\Omega_{\mathfrak{m}}) = \dim(\Lambda_{\mathfrak{m}})\)，于是对 b) 取上确界即得（§ 3, 第 2 节，命题 3）。

命题 2.—— 设 A 为 Noether 环，Ω 为对偶化 A-模。对 A 的每个素理想 p，A\(_{p}\)-模 Ω\(_{p}\) 是对偶化的。

考虑 A 的素理想的一条饱和链 \(\mathfrak{p} \subset \mathfrak{p}_1 \subset \ldots \subset \mathfrak{p}_r\)，使得理想 \(\mathfrak{p}_r\) 是极大的。对 \(r\) 作归纳，可设 \(\mathrm{A}_{\mathfrak{p}_1}\)-模 \(\Omega_{\mathfrak{p}_1}\) 是对偶化的。把 \(\Lambda\) 换成 \(\mathrm{A}_{\mathfrak{p}_1}\)，把 \(\mathfrak{p}\) 换成 \(\mathfrak{p}\mathrm{A}_{\mathfrak{p}_1}\)，就化归为环 A 是局部环且链 \(\mathfrak{p} \subset \mathfrak{m}_{\Lambda}\) 为饱和链的情形。

于是设 \(d = \dim(A) = \mathrm{ht}(\mathfrak{m}_{\mathbf{A}})\)。由于 A 是 Macaulay 环（§ 2, 第 2 节，命题 2 的推论），有 \(\dim(A_{\mathfrak{p}}) = \mathrm{ht}(\mathfrak{p}) = d - 1\)。对每个整数 \(i\)，\(A_{\mathfrak{p}}\)-模 \(\operatorname{Ext}_{A_{\mathfrak{p}}}^{i}(\kappa(\mathfrak{p}), \Omega_{\mathfrak{p}})\) 同构于 \(\operatorname{Ext}_{\Lambda}^{i}(A/\mathfrak{p}, \Omega)_{\mathfrak{p}}\)。于是由 § 3, 第 2 节，命题 2，只需证明 \(A/\mathfrak{p}\)-模 \(\operatorname{Ext}_{A}^{i}(A/\mathfrak{p}, \Omega)\) 在 \(i \neq d - 1\) 时为零，而在 \(i = d - 1\) 时秩为 1。

设 \(x\) 为 \(\mathfrak{m}_{\mathrm{A}} - \mathfrak{p}\) 的一个元素，并用 \(\overline{x}\) 表示它在 \(A / \mathfrak{p}\) 中的类。考虑 A-模的正合列

\[
0 \rightarrow \mathrm{A} / \mathfrak {p} \xrightarrow {\overline {{{x}}}} \mathrm{A} / \mathfrak {p} \longrightarrow \mathrm{A} / (\mathfrak {p} + x \mathrm{A}) \rightarrow 0.
\]

A-模 \(\mathrm{A} / (\mathfrak{p} + x\mathrm{A})\) 是有限长度的，因为它的支撑集化为 \(\mathfrak{m}_{\mathrm{A}}\)；由于 A-模 \(\Omega\) 是对偶化的，有 \(\operatorname{Ext}_{\mathrm{A}}^{i}(\mathrm{A} / (\mathfrak{p} + x\mathrm{A}), \Omega) = 0\)（当 \(i \neq d\) 时）（§ 8, 第 5 节，例 3）。于是，由与上述序列以及与 \(\Omega\) 相伴的扩张模正合列推出：以 \(x\) 为比的乘法映射在 A-模 \(\operatorname{Ext}_{\mathrm{A}}^{i}(\mathrm{A} / \mathfrak{p}, \Omega)\) 上是满射（当 \(i \neq d - 1\) 时），这蕴含该模为零（Nakayama 引理）。特别地，\(\operatorname{Ext}_{\mathrm{A}}^{d}(\mathrm{A} / \mathfrak{p}, \Omega)\) 为零，从而得到正合列

\[
0 \rightarrow \operatorname{Ext} _ {\Lambda} ^ {d - 1} (\mathrm{A} / \mathfrak {p}, \Omega) \xrightarrow {x} \operatorname{Ext} _ {\Lambda} ^ {d - 1} (\mathrm{A} / \mathfrak {p}, \Omega) \longrightarrow \operatorname{Ext} _ {\Lambda} ^ {d} (\mathrm{A} / (\mathfrak {p} + x \Lambda), \Omega) \rightarrow 0.
\]

我们有 \(\operatorname{long}_{\mathrm{A}}(\operatorname{Ext}_{\mathrm{A}}^{i}(\mathrm{A}/(\mathfrak{p}+x\mathrm{A}),\Omega))=\operatorname{long}_{\mathrm{A}}(\mathrm{A}/(\mathfrak{p}+x\mathrm{A}))\)（前引）；于是命题由下面的引理得出，把该引理应用于环 B = A/p 及 B-模 \(M=\operatorname{Ext}_{\Lambda}^{d-1}(\mathrm{A}/\mathfrak{p},\Omega)\)：

引理 1. 设 B 为一维的局部 Noether 整环，M 为有限生成的无扭 B-模。假设对 B 的每个非零元素 x，都有 long \(_{B}\) (M/xM) = long \(_{B}\) (B/xB)。那么 B-模 M 的秩为 1。

事实上，设 \(r\) 为 M 的秩；存在 M 的一个秩为 \(r\) 的自由子模 L，使 M/L 为扭模（VII, § 4, 第 1 节，命题 1 的推论），从而是有限长度的（VII, § 2, 第 5 节，引理 1）。M/L 的零化子不化为 0，因而含有 \(x\) 的一个非零元素 \(\mathfrak{m}_{\mathbb{R}}\)。考虑图表

交换图表

\[
\begin{array}{c c c c c c c} 0 \to & \mathrm{L} & \longrightarrow & \mathrm{M} & \longrightarrow & \mathrm{M} / \mathrm{L} & \to 0 \\ & \Biggl \downarrow x _ {\mathrm{L}} & & \Biggl \downarrow x _ {\mathrm{M}} & & \Biggl \downarrow 0 \\ 0 \to & \mathrm{L} & \longrightarrow & \mathrm{M} & \longrightarrow & \mathrm{M} / \mathrm{L} & \to 0. \end{array}
\]

由蛇形引理（A, X, 第 4 页，命题 2），可导出一个正合列

\[
0 \rightarrow \mathrm{M} / \mathrm{L} \longrightarrow \mathrm{L} / x \mathrm{L} \longrightarrow \mathrm{M} / x \mathrm{M} \longrightarrow \mathrm{M} / \mathrm{L} \rightarrow 0,
\]

由此得 long(M/xM) = long(L/xL)。由于按假设 long(M/xM) = long(B/xB)，而 long(L/xL) = r·long(B/xB)，故得 r = 1。

推论 1. 对 A 的每个乘性子集 S，\(\mathrm{S}^{-1}\mathrm{A}\)-模 \(\mathrm{S}^{-1}\Omega\) 都是对偶化的。

推论 2.—— Ω 的支撑集等于 Spec(A)。

事实上，按定义，局部环上的对偶化模是非零的。

推论 3. 设 M 为有限生成的 A-模，i 为整数。A-模 \(\operatorname{Ext}_{\mathrm{A}}^{i}(\mathrm{M},\Omega)\) 是有限生成的，且其支撑集的余维数 \(\geqslant i\)（在 \(\operatorname{Spec}(\mathrm{A})\) 中）。

第一个断言由 A, X, 第 108 页，推论可得。设 \(\mathfrak{p}\) 为 \(\operatorname{Ext}_{\mathrm{A}}^{i}(\mathrm{M},\Omega)\) 的支撑集的一个素理想。有 \(\operatorname{Ext}_{\mathrm{A}}^{i}(\mathrm{M},\Omega)_{\mathfrak{p}} \neq 0\)，从而 \(\operatorname{Ext}_{\mathrm{A}_{\mathfrak{p}}}^{i}(\mathrm{M}_{\mathfrak{p}},\Omega_{\mathfrak{p}}) \neq 0\)（\(\S 3, \mathrm{n}^{\circ}2\)，命题 2），这蕴含 \(\mathrm{di}_{\mathrm{A}_{\mathfrak{p}}}(\Omega_{\mathfrak{p}}) \geqslant i\)。由于 \(\Omega_{\mathfrak{p}}\) 是 \(\mathrm{A}_{\mathfrak{p}}\)-模且是对偶化的（命题 2），有 \(\mathrm{di}_{\mathrm{A}_{\mathfrak{p}}}(\Omega_{\mathfrak{p}}) = \dim(\mathrm{A}_{\mathfrak{p}})\)（命题 1），由此得推论。

命题 3.—— 设 A 为 Noether 局部环，Ω 为对偶化 A-模，而 M 为有限生成的 A-模。

\[
\text { a) } O n a \operatorname{Ext} _ {\mathrm{A}} ^ {i} (\mathrm{M}, \Omega) = 0 \text { for } i <   \dim (\mathrm{A}) \dots \dim_ {\mathrm{A}} (\mathrm{M}).
\]

\begin{enumerate}
\def\labelenumi{\alph{enumi})}
\setcounter{enumi}{1}
\tightlist
\item
  令 \(c = \dim(\mathrm{A}) - \dim_{\mathrm{A}}(\mathrm{M})\)。若 M 非零，则 A-模 \(\operatorname{Ext}_{\mathrm{A}}^{c}(\mathrm{M}, \Omega)\) 不为零。
\end{enumerate}

\[
\mathrm{c)} O n a \operatorname{Ext} _ {\mathrm{A}} ^ {i} (\mathrm{M}, \Omega) = 0 p o u r i > \dim (\mathrm{A}) - \operatorname{prof} _ {\mathrm{A}} (\mathrm{M}).
\]

设 M 非零，并用 F 表示它的支撑集。由 § 1, 第 5 节，命题 9，断言 a) 与 b) 的合取等价于 \(\operatorname{prof}_{\mathbb{F}}(\Omega)=c\)。由于 \(\Omega\) 是 Macaulay 模且其支撑集等于 \(\operatorname{Spec}(\mathbb{A})\)（命题 1 及命题 2 的推论 2），我们有

\[
\operatorname{prof} _ {\mathrm{F}} (\Omega) = \operatorname{codim} (\mathrm{F}, \operatorname{Spec} (\Lambda)) = c
\]

（§ 2, 第 1 节，命题 1 的推论及第 2 节，命题 2 的推论）。

我们对 M 的深度作归纳来证明 c)。若 \(\mathrm{prof}_{\Lambda}(\mathbf{M}) = 0\)，则确实有 \(\mathrm{Ext}_{\Lambda}^{i}(\mathbf{M},\Omega) = 0\)（当 \(i > \dim (\mathbf{A})\) 时），因为 \(\mathrm{di}_{\Lambda}(\Omega) = \dim (\mathbf{A})\)（命题 1）。设 \(\mathrm{prof}_{\Lambda}(\mathbf{M}) > 0\)，则存在 \(x\) 的一个元素 \(\mathfrak{m}_{\mathbf{A}}\)，使以 \(x\) 为比的乘法映射在 M 中是单射的。我们有 \(\mathrm{prof}_{\Lambda}(\mathbf{M} / x\mathbf{M}) = \mathrm{prof}_{\Lambda}(\mathbf{M}) - 1\)（\(\S 1, \mathrm{n}^{\circ}4\)，命题 7）。

考虑扩张模的正合列

\[
\operatorname{Ext} _ {\mathrm{A}} ^ {i} (\mathrm{M}, \Omega) \xrightarrow {x} \operatorname{Ext} _ {\mathrm{A}} ^ {i} (\mathrm{M}, \Omega) \longrightarrow \operatorname{Ext} _ {\Lambda} ^ {i + 1} (\mathrm{M} / x \mathrm{M}, \Omega)
\]

它相伴于正合列

\[
0 \to \mathrm{M} \xrightarrow {x} \mathrm{M} \longrightarrow \mathrm{M} / x \mathrm{M} \to 0.
\]

当 \(i > \dim(\mathrm{A}) - \operatorname{prof}_{\mathrm{A}}(\mathrm{M})\) 时，A-模 \(\operatorname{Ext}_{\mathrm{A}}^{i+1}(\mathrm{M}/x\mathrm{M},\Omega)\) 按归纳假设为零，于是以 \(x\) 为比的乘法映射在 \(\operatorname{Ext}_{\mathrm{A}}^{i}(\mathrm{M},\Omega)\) 上是满射的，这蕴含该 A-模为零（Nakayama 引理）。这就证明了 c)。

推论.—— 若 M 是 Macaulay 模，则当 i ≠ c 时有 Ext \(^{i}\) A(M, Ω) = 0；A-模 Ext \(^{c}\) A(M, Ω) 是 Macaulay 模，且其支撑集等于 M 的支撑集。

第一个断言由命题 3, a) 与 c) 得出。设 \(\mathfrak{p} \in \operatorname{Supp}(M)\)；把 § 2, 第 1 节，命题 1 应用于 M 与 A，我们有

\[
\dim (A _ {p}) - \dim_ {A _ {p}} (M _ {p}) = \dim (A) - \dim_ {A} (M) = c;
\]

由于 \(A_{p}\)-模 \(\Omega_{p}\) 是对偶化的（命题 2），由命题 3, b) 可知 \(A_{p}\)-模 \(\operatorname{Ext}_{A_{p}}^{c}(M_{p},\Omega_{p})\) 不为零。因此，\(\operatorname{Ext}_{A}^{c}(M,\Omega)\) 的支撑集等于 M 的支撑集。

最后，我们对 \(\dim(M)\) 作归纳，证明 A-模 \(\operatorname{Ext}_{A}^{c}(M,\Omega)\) 是 Macaulay 模。当 \(\dim(M)=0\) 时断言成立，因为每个有限长度模都是 Macaulay 模。设 \(\dim(M)>0\)，并选取 \(x\) 的一个元素 \(\mathfrak{m}_{\mathrm{A}}\)，使乘法映射 \(x_{\mathrm{M}}\) 是单射的。A-模 M/\(x\mathrm{M}\) 是 Macaulay 模（\(\S 2\) 第 3 节，命题 4），其维数为 \(\dim(M)-1\)。考虑到上述结果，与正合列 \(0\to\mathrm{M}\xrightarrow{x}\mathrm{M}\longrightarrow\mathrm{M}/x\mathrm{M}\to0\) 相伴的扩张模正合列化为

\[
0 \to \operatorname{Ext} _ {\mathrm{A}} ^ {c} (\mathrm{M}, \Omega) \xrightarrow {x} \operatorname{Ext} _ {\mathrm{A}} ^ {c} (\mathrm{M}, \Omega) \longrightarrow \operatorname{Ext} _ {\mathrm{A}} ^ {c + 1} (\mathrm{M} / x \mathrm{M}, \Omega) \to 0;
\]

于是 § 2, 第 3 节的命题 4 与归纳假设蕴含 Ext \(_{A}^{c}\) (M, Ω) 是 Macaulay 模，由此得推论。

